\documentclass[10pt]{article}
\usepackage[margin=1in]{geometry}
\usepackage{booktabs}
\usepackage{longtable}
\usepackage{hyperref}
\usepackage{xcolor}
\usepackage{parskip}

\newcommand{\repl}{{\color{green!50!black}\textbf{REPRODUCED}}}
\newcommand{\part}{{\color{orange!80!black}\textbf{PARTIAL}}}
\newcommand{\contra}{{\color{red!80!black}\textbf{CONTRADICTED}}}
\newcommand{\qual}{{\color{blue!70!black}\textbf{QUALITATIVELY REPRODUCED}}}
\newcommand{\nope}{{\color{gray}\textbf{NOT ATTEMPTED}}}

\title{BVBRC-123 Replication Report:\\ \textit{Aeromonas veronii} Alim\_AV\_1000 (Rahman et al., 2023)}
\author{X-100 Replication Project, Wave 2026-07-05 (BVBRC set)}
\date{2026-07-05}

\begin{document}
\maketitle

\begin{abstract}
Independent replication of the whole-genome sequence analysis reported by Rahman et al.\ (\textit{J Adv Vet Anim Res} 2023, PMID 37969805, PMC10636080). Assembly-level and coarse annotation claims reproduce exactly from the deposited genome (GCA\_026738955.1). \textbf{Verdict: PARTIAL.} Two claims are contradicted at the genomic level: (a) the MLST ST 492 designation does not match any allele of the deposited assembly's actual profile, and (b) the paper's stated BioProject accession PRJNA810265 points to a different organism (\textit{Pasteurella multocida} DC2020 by the same institution). Prophage and virulence claims are qualitatively reproduced. Wet-lab pathogenicity claims out of scope.
\end{abstract}

\section{Paper Under Test}
\textbf{Rahman MM, Sadekuzzaman M, Rahman MA, Siddique MP, Uddin MA, Haque ME, Chowdhury MGA, Khasruzzaman AKM, Rahman MT, Hossain MT, Islam MA.} ``Complete genome sequence analysis of the multidrug resistant \textit{Aeromonas veronii} isolated for the first time from stinging catfish (Shing fish) in Bangladesh.'' \textit{J Adv Vet Anim Res} 10(3): 570--578 (Sep 2023). DOI: \href{https://doi.org/10.5455/javar.2023.j711}{10.5455/javar.2023.j711}. PMID 37969805.

\section{Claim-by-Claim Verdict Table}

\begin{longtable}{@{}p{0.5cm}p{7cm}p{4cm}p{4cm}@{}}
\toprule
\# & Claim & My finding & Verdict \\
\midrule
\endhead

C1 & Genome size 4{,}494{,}515 bp & 4{,}494{,}464 bp (Δ = −51 bp, 0.001\%) & \repl \\
C2 & GC content 58.87\% & 58.87\% & \repl \\
C3 & 93 contigs & 93 & \repl \\
C4 & N50 = 150{,}337 & 150{,}337 & \repl \\
C5 & L50 = 12 & 12 & \repl \\
C6 & 0 plasmids & 0 (PlasmidFinder) & \repl \\
C7 & 4{,}229 CDS (RAST) & 4{,}099 CDS (PGAP); tool variance ~3\% & \part \\
C8 & 102 tRNA & 102 & \repl \\
C9 & 13 rRNA & 28 (PGAP splits per-contig fragments) & \part \\
C10 & MLST ST 492 & Zero of 6 ST 492 alleles match deposited genome; scan of all 2{,}756 aeromonas profiles finds none with $\geq$3 of my alleles & \contra \\
C11 & MDR phenotype linked to acquired resistance genes & Acquired: OXA-12, cphA4 ($\beta$-lactam/carbapenem only). No acquired tet/qnr/aminoglycoside genes above threshold. Paper's higher CARD count likely includes chromosomal targets. & \part \\
C12 & Closest to TH0426 + B56 (paper name; likely B565) & TH0426 ANI 96.34\%, B565 96.34\%, FDAARGOS\_632 96.47\% --- all essentially equidistant & \part \\
C13 & 2 intact + 1 incomplete prophage (PHASTER) & 2 contigs with dense phage-gene clusters (16 + 7 genes), several 3-gene fragments & \qual \\
C14 & BioProject PRJNA810265; BioSample SUB11126221 & PRJNA810265 = \textit{Pasteurella multocida} DC2020 (WRONG); real BioProject = PRJNA827572; real BioSample = SAMN27611687 & \contra \\
C15 & 100\% IP mortality, 80--100\% oral, 0\% control in fish infection & Wet-lab, requires live isolate + BSL-2 aquarium & \nope \\
C16 & Proteome $\geq$95\% conserved vs NZ\_CP044060.1 & Not attempted (OrthoFinder deferred) & \nope \\

\bottomrule
\end{longtable}

\section{Method Summary}
All computations on CherryRd (macOS Homebrew): \texttt{abricate 1.4.0} (DBs dated 2026-07-03), \texttt{skani 0.3.2}, \texttt{fastANI 1.34}, \texttt{pdftotext} (Poppler), Python 3.14 stdlib for FASTA parsing. Assembly + PGAP annotation downloaded via NCBI Datasets v2 REST. MLST via pubMLST REST (POST base64-encoded FASTA to \texttt{schemes/1/sequence} endpoint). Phylogenetic comparators (TH0426, B565, FDAARGOS\_632) downloaded via NCBI Datasets v2. LLM-judge via \texttt{argo:gpt-5.4} on cherryrd litellm aggregator (\texttt{http://<tailnet-aggregator>:4000/v1}). Zero paid API calls.

\section{Numerical Results (Assembly)}
Independent recomputation from the downloaded RefSeq FASTA \texttt{GCF\_026738955.1\_ASM2673895v1\_genomic.fna}:

\begin{center}
\begin{tabular}{@{}lrr@{}}
\toprule
Metric & Paper & This work \\
\midrule
Contigs & 93 & 93 \\
Total bp & 4{,}494{,}515 & 4{,}494{,}464 \\
Largest contig & (not stated) & 296{,}612 \\
Smallest contig & (not stated) & 208 \\
N50 & 150{,}337 & 150{,}337 \\
L50 & 12 & 12 \\
GC \% & 58.87 & 58.87 \\
Coding density \% & (not stated) & 88.17 \\
CheckM completeness \% & (not stated) & 98.96 \\
CheckM contamination \% & (not stated) & 1.38 \\
\bottomrule
\end{tabular}
\end{center}

\section{Numerical Results (MLST --- the CONTRADICTION)}

pubMLST scheme 1 (Aeromonas MLST, 6-locus: gyrB, groL, gltA, metG, ppsA, recA) applied to the deposited assembly:

\begin{center}
\begin{tabular}{@{}lll@{}}
\toprule
Locus & My exact allele match & ST 492 expected allele \\
\midrule
gyrB & 633 & 112 \\
groL & 91  & 347 \\
gltA & 340 & 44  \\
metG & 124 & 217 \\
ppsA & (no exact match) & 384 \\
recA & 1460 & 381 \\
\bottomrule
\end{tabular}
\end{center}

Zero shared alleles. Additionally, scanning all 2{,}756 profiles in the pubMLST aeromonas database returned zero profiles with $\geq$3 of my alleles --- this appears to be an unregistered/novel MLST profile, not merely a mis-typed variant of ST 492. Either the paper reports the wrong ST or a different physical isolate was deposited.

\section{Numerical Results (ANI)}

\begin{center}
\begin{tabular}{@{}lrr@{}}
\toprule
Reference & skani ANI (\%) & fastANI (\%) \\
\midrule
A. veronii TH0426 (GCF\_001593245.1) & 96.34 & 96.24 \\
A. veronii B565 (GCF\_000204115.1) & 96.34 & 96.34 \\
A. veronii FDAARGOS\_632 (GCF\_008693705.1) & 96.47 & 96.38 \\
A. veronii bv. veronii type (GCA\_001908535.1, via NCBI ANI check) & --- & 96.30 \\
\bottomrule
\end{tabular}
\end{center}

All comparators sit within a 0.14\% ANI window. TH0426 is not uniquely closest; paper's phylogenetic-tree-based closest-sibling claim reflects branch topology, not distinctive genome-distance.

\section{What Worked, What Didn't}

\textbf{Worked first-shot:}
Europe PMC PDF fetch; pdftotext extraction; NCBI Datasets v2 REST downloads (all 5 genomes); Python assembly-stats recompute (bit-exact N50, L50, GC\%); abricate on 5 databases; pubMLST REST (after switching to JSON body + base64 encoding); skani and fastANI; LLM-judge via aggregator gpt-5.4.

\textbf{Failed and worked-around:}
NCBI PMC direct PDF (returned HTML); OpenClaw \texttt{pdf} tool (paid providers all unavailable); local \texttt{mlst} binary (Homebrew Perl XS mismatch); pubMLST REST form-encoded body (rejected); local Argo :44497 and cherryrd claude-opus-4.8 (both 502-ing during this run).

\textbf{Not attempted (documented as extensions):}
Prokka/Bakta reannotation on same assembly to isolate annotator variance for CDS count; PHASTER/PHASTEST full rerun for exact prophage coordinates; full RaxML tree with PATRIC PGFams to validate closest-sibling topology; OrthoFinder for proteome-conservation claim; wet-lab fish infection.

\section{Open Questions}
See \texttt{report/open\_questions.json} for structured JSON list. Summary:
\begin{enumerate}
\item Did a sample mix-up occur between MLST typing and NCBI deposition? (Basis: ST 492 alleles share nothing with deposited genome.)
\item How pervasive are cross-project accession errors in single-institution WGS papers? (Basis: PRJNA810265 points to wrong organism.)
\item How well does PATRIC k-mer AMR predict phenotype when acquired mobile AMR genes are absent? (Basis: MDR phenotype with only $\beta$-lactam acquired-gene support.)
\item What phylogenetic method reliably resolves closest-relative when ANI is tied to $\sim$0.15\%? (Basis: TH0426, B565, FDAARGOS\_632 all essentially equidistant.)
\item Would long-read resequencing coalesce fragmented prophage regions in this contig-level assembly? (Basis: 2 dense + 6 sparse phage-gene contigs, 93 contigs suggest repeat-driven fragmentation.)
\end{enumerate}

\end{document}
